\section{国际海上货物运输法律制度}
\begin{enumerate}
  \item 货物运输法要求的多一些。海运法、国际铁路货物运输法、国际航空货物运输法。
  \item 运输方式。
        \begin{enumerate}
          \item 班轮运输。类似地铁。固定的始发港目的港。小中货间杂货。
          \item 租船运输。打出租。航次跑一个航程，定期跑一段时间。租船合同。
        \end{enumerate}
  \item 单证类型。\zd{名字含“提单”才是物权凭证}（内容上记载所有权归属）。
        \begin{kk}{海运提单}
          用以证明海上货物运输\zd{合同的订立}和货物已经由承运人接管 \hz 装船以及承运人保证据以\zd{交付货物}的单证。
          \begin{enumerate}
            \item 是海上\zd{运输合同的证明}。
            \item 是承运人出具的\zd{接收货物的收据}。
            \item 是承运人\zd{交付货物的物权凭证}。见正本提单放货。\zd{转让海运提单 = 转让在途运输的货物所有权}。
          \end{enumerate}
          \tcblower
          分类标准。
          \begin{enumerate}
            \item 是否已实际装船。\zd{已装船提单}（船名和航次），\zd{收货待运提单/备运提单}（船名、航次空着）。
            \item 有无不良批注。\zd{清洁提单和不清洁提单}。\al{少一箱，部分纸箱破损}。
            \item 按收货人抬头。\zd{记名/不记名/指示。收货人/空白 \hz 交持有人/凭指示。不能/交付/背书转让}。
            \item \zd{倒签提单和预借提单}。装船时间作假。已/未装船。往前倒着签/借一个。用的少。考过一次。
          \end{enumerate}
        \end{kk}

        \begin{kk}{海运单}
          航程较短的运输。单货一起。\zd{不是物权凭证。不能单方转让}。都具有运输合同证明和货物收据的作用。
        \end{kk}
  \item 承运人责任。\zd{违反最低限度的义务 \yjt 有无过失 \yjt 航行过失 \yjt 火灾过失 \yjt 看规则}。
        \begin{enumerate}
          \item \zd{交货责任}。货交给谁、交错了怎么办、赔不赔、怎么赔。
                \begin{enumerate}
                  \item \zd{无单放货}。\zd{无正本提单}交付货物案件。
                        \begin{enumerate}
                          \item 与无单取货人承担\zd{连带责任}。违约和侵权竞合。实践中一般选择告承运人（方便）。
                          \item 赔偿额以 CIF 价计算\zd{足额}赔偿（\zd{不用海事赔偿责任限制，防止恶意串通}）。
                        \end{enumerate}
                  \item 免责。
                        \begin{enumerate}
                          \item 依法行为/依法免责。\zd{被}海关、法院、商检机关等\zd{依法处理。过时效无人清关}。
                          \item \zd{执行记名提单}托运人\zd{指示}终止运输、返还货物、变更到达地、交其他收货人。
                          \item \zd{已有正本提单已提货}。一式数份正本提单（会计入账需求）。只认第一份。
                        \end{enumerate}
                  \item 诉讼时效 \zd{1} 年。\zd{适用我国海商法 \yjt 其他法律}。
                \end{enumerate}
          \item \zd{货损责任}。在运输途中出现了毁损灭失，赔不赔，怎么赔。海牙规则、维斯比规则、汉堡规则。
                \begin{kk}{海牙规则/维斯比规则}
                  \begin{enumerate}
                    \item \zd{所有海运规则都包含}（包括我国海商法）。
                    \item 最低限度的义务/强制性规定。\zd{适航义务，管货义务}。\zd{这两项义务是强制性的}。
                    \item 责任期间。装到卸。
                    \item 免责。
                          \begin{enumerate}
                            \item 共有 17 项，2 项过失（\zd{航行过失、火灾过失}）免责，15 项无过失免责。因为便宜。
                            \item \zd{船上自然人是承运人的雇佣人、代理人}。船公司、远洋运输公司才是承运人。
                          \end{enumerate}
                    \item 火灾责任细分 4 种。
                          \begin{enumerate}
                            \item 遭\zd{雷击失火不赔}；雇佣人、代理人\zd{过失不赔}。\al{烟头引发}。\xjt 都赔。
                            \item \zd{承运人实际过失}导致的。解决实践中\al{承运人提供的船舶有问题}。
                            \item \zd{谋私}。\al{故意纵火}。
                          \end{enumerate}
                  \end{enumerate}
                  \tcblower
                  \begin{enumerate}
                    \item 未对迟延、实际承运人、舱面货和活牲畜、保函进行规定。个案分析。
                    \item 诉讼时效 \zd{1} 年。
                    \item 维斯比规则是对海牙规则的非实质性修改。责任限制及最高赔偿限额略高。
                  \end{enumerate}
                \end{kk}

                \begin{kk}{汉堡规则}
                  \begin{enumerate}
                    \item 去掉了过失免责。责任期间为接到交。
                    \item 对迟延责任、舱面货和活牲畜、\zd{保函（可考性最强）}进行了\zd{强制性规定}。有约定则约定无效。
                    \item 迟延交付的赔偿为\zd{迟交货物运费的 2.5 倍}，但不应超过应付运费的总额。
                    \item \zd{实际承运人与承运人共负连带责任}。
                    \item 诉讼时效 \zd{2} 年。
                  \end{enumerate}

                  \begin{kk}{承认善意保函}
                    \al{神华煤 10W 吨，船老板经验认为只有 8 吨，神华说就有 10W 吨，我给你开保函，你别动笔。这里没有恶意串通，这就是双方发生分歧}。
                    \begin{enumerate}
                      \item \zd{以善意保函换取清洁提单}。
                      \item 仅在\zd{托运人和承运人}之间有效，\zd{不能对抗收货人}。\zd{合同相对性原理}。
                      \item 两段法律关系。承收之间是提单，托承之间才是保函。
                    \end{enumerate}
                  \end{kk}
                \end{kk}

        \end{enumerate}
\end{enumerate}
